\documentclass[12pt]{article}
\usepackage[utf8]{inputenc}
\usepackage[T1]{fontenc}
\usepackage{hyperref}
\title{Mendalam Kayan Language Resources for K-12 Curriculum Development}
\author{Compiled by Claude Sonnet 5.5 (claude-sonnet-5-5) for Andr\'as Kornai}
\date{}
\begin{document}
\maketitle

\section{Introduction}
Mendalam Kayan (ISO 639-3 xkd; Glottolog mend1267) is an Austronesian language of the Kayan-Murik group. It is spoken in the Mendalam area of Kapuas Hulu regency, West Kalimantan, Indonesia. The only speaker figure I found (35,000) covers the whole Kayan cluster, not Mendalam alone. I found no school materials in it.

\section{Primary Educational Resources}
I found no primers, textbooks or school programmes for Mendalam Kayan.

\section{Linguistic and Reference Materials}
\begin{itemize}
\item Wikipedia on the Kayan language cluster (Borneo), which lists xkd among its varieties. \url{https://wikipedia.com/wiki/ISO_639:bfg}
\item Wikipedia, ``Languages of Kalimantan'': regional context. \url{https://en.wikipedia.org/wiki/Languages_of_Kalimantan}
\end{itemize}

\section{Religious and Community Resources}
\begin{itemize}
\item Joshua Project profile of the Mendalam Kayan. It reports that no Christian materials exist in the language. It is a missions database, so read it with that framing in mind. \url{https://joshuaproject.net/people_groups/12592}
\item Kajo Magazine, a feature on a Kayan encounter at Putussibau: cultural background only. \url{https://kajomag.com/kayan-putussibau-indonesia/}
\end{itemize}

\section{Digital and Audio Resources}
\begin{itemize}
\item CREM-CNRS archive: 1997 Kayan song recordings made in Tanjung Durian, Mendalam. The archive's metadata, not its content, is what I saw. Check the access terms. \url{https://archives.crem-cnrs.fr/archives/items/CNRSMH_I_2011_016_006_06/}
\end{itemize}

\section{Curriculum Development Context}
Very little is documented for this variety. Related Kayan varieties may offer descriptions to adapt, but their forms differ from Mendalam. A curriculum would need community-led documentation and the involvement of local speakers.

\end{document}